% 內科部臨床技術分類表格
% 需要的套件：longtable, booktabs, array

\documentclass[a4paper,12pt]{article}
\usepackage[utf8]{inputenc}
\usepackage{xeCJK}
\usepackage{longtable}
\usepackage{booktabs}
\usepackage{array}
\usepackage{geometry}
\usepackage{caption}

% 設定中文字型
\IfFontExistsTF{Noto Serif CJK TC}{
  \setCJKmainfont{Noto Serif CJK TC}
}{\setCJKmainfont{SimSun}}

% 頁面設定
\geometry{
    a4paper,
    left=2cm,
    right=2cm,
    top=2cm,
    bottom=2cm
}

\begin{document}

\section*{內科部侵入性臨床技術檢核表}
\subsection*{依技能類別分類整理}

\textbf{監督等級說明：}
\begin{itemize}
\item \textbf{1} = Independent operation（獨立操作）
\item \textbf{2} = Supervisor available（主治醫師可聯繫）
\item \textbf{3} = Supervisor direct supervision（主治醫師直接監督）
\item \textbf{4} = Observer（觀察者）
\end{itemize}

\vspace{1em}

% 完整技能表格
\begin{longtable}{>{\raggedright}p{3cm}>{\raggedright\arraybackslash}p{4.5cm}ccccccc}
\caption{內科部臨床技術監督等級表} \label{tab:internal-medicine-skills} \\
\toprule
\textbf{技能類別} & \textbf{技能名稱} & \textbf{M5} & \textbf{M6} & \textbf{PGY1} & \textbf{PGY2} & \textbf{R1} & \textbf{R2} & \textbf{R3} \\
\midrule
\endfirsthead

\multicolumn{9}{c}{\textit{續前頁}} \\
\toprule
\textbf{技能類別} & \textbf{技能名稱} & \textbf{M5} & \textbf{M6} & \textbf{PGY1} & \textbf{PGY2} & \textbf{R1} & \textbf{R2} & \textbf{R3} \\
\midrule
\endhead

\midrule
\multicolumn{9}{r}{\textit{續下頁}} \\
\endfoot

\bottomrule
\endlastfoot

% 基本照護技術
\textbf{基本照護技術} & & & & & & & & \\
\midrule
& 導尿管置放術 \newline \textit{Foley insertion} & 4 & 2 & 1 & 1 & 1 & 1 & 1 \\
& 鼻胃管置放術 \newline \textit{NG tube} & 4 & 2 & 1 & 1 & 1 & 1 & 1 \\
& 抽血 \newline \textit{Blood Drawing} & 4 & 2 & 1 & 1 & 1 & 1 & 1 \\
\midrule

% 血管通路技術
\textbf{血管通路技術} & & & & & & & & \\
\midrule
& 中心靜脈導管放置 \newline \textit{CVC placement} & 4 & 4 & 4 & 3 & 2 & 1 & 1 \\
& 雙腔靜脈導管 \newline \textit{Double lumen catheter} & 4 & 4 & 4 & 4 & 3 & 2 & 1 \\
& 動脈出血控制 \newline \textit{Control of bleeding (Arterial)} & 4 & 4 & 3 & 3 & 2 & 1 & 1 \\
\midrule

% 急救技術
\textbf{急救技術} & & & & & & & & \\
\midrule
& 心肺腦復甦術 \newline \textit{CPCR} & 4 & 4 & 3 & 3 & 2 & 2 & 1 \\
& 氣管內管更換 \newline \textit{Endotracheal tube change} & 4 & 4 & 4 & 4 & 3 & 2 & 1 \\
& 胃食道球置入 \newline \textit{S-B tube} & 4 & 4 & 4 & 4 & 3 & 2 & 1 \\
\midrule

% 診斷性穿刺技術
\textbf{診斷性穿刺技術} & & & & & & & & \\
\midrule
& 腹水引流術 \newline \textit{Paracentesis} & 4 & 4 & 3 & 3 & 2 & 1 & 1 \\
& 胸腔穿刺術 \newline \textit{Thoracentesis} & 4 & 4 & 4 & 3 & 3 & 2 & 1 \\
& 關節腔液抽吸 \newline \textit{Synovial fluid aspiration} & 4 & 4 & 4 & 4 & 4 & 3 & 2 \\
& 腰椎穿刺 \newline \textit{Lumbar puncture} & 4 & 4 & 4 & 4 & 3 & 2 & 1 \\

\end{longtable}

\vspace{1em}

\section*{技能學習進程摘要}

\subsection*{各訓練階段可獨立操作技能統計}

\begin{table}[h]
\centering
\begin{tabular}{lc}
\toprule
\textbf{訓練階段} & \textbf{可獨立操作技能數} \\
\midrule
M5 & 0項 \\
M6 & 0項 \\
PGY1 & 3項 \\
PGY2 & 3項 \\
R1 & 3項 \\
R2 & 6項 \\
R3 & 12項 \\
\bottomrule
\end{tabular}
\caption{各階段可獨立操作技能統計}
\end{table}

\subsection*{技能類別統計}

\begin{table}[h]
\centering
\begin{tabular}{lc}
\toprule
\textbf{技能類別} & \textbf{技能數量} \\
\midrule
基本照護技術 & 3項 \\
血管通路技術 & 3項 \\
急救技術 & 3項 \\
診斷性穿刺技術 & 4項 \\
\midrule
\textbf{總計} & \textbf{13項} \\
\bottomrule
\end{tabular}
\caption{技能類別統計表}
\end{table}

\subsection*{技能難度分層}

\subsubsection*{初階技能（PGY1即可獨立操作）}
\begin{enumerate}
\item 導尿管置放術 (Foley insertion)
\item 鼻胃管置放術 (NG tube)
\item 抽血 (Blood Drawing)
\end{enumerate}

\subsubsection*{中階技能（R2可獨立操作）}
\begin{enumerate}
\item 中心靜脈導管放置 (CVC placement)
\item 動脈出血控制 (Control of bleeding - Arterial)
\item 腹水引流術 (Paracentesis)
\end{enumerate}

\subsubsection*{進階技能（R3才可獨立操作）}
\begin{enumerate}
\item 雙腔靜脈導管 (Double lumen catheter)
\item 心肺腦復甦術 (CPCR)
\item 氣管內管更換 (Endotracheal tube change)
\item 胃食道球置入 (S-B tube)
\item 胸腔穿刺術 (Thoracentesis)
\item 腰椎穿刺 (Lumbar puncture)
\end{enumerate}

\vspace{1em}
\noindent\textit{資料來源：醫療科部Invasive skills checklist教學部彙整（修訂日期：2022/12/13）}

\end{document}